% !PRINTED none

\begin{center}
\gesperrt{Siebzehnter Abschnitt.}

\gesperrt{Theorie der algebraischen K\"orper.}

\rule{4cm}{0.4pt}

\S.\,144.

\gesperrt{Basis eines algebraischen Zahlk\"orpers.}

\gesperrt{Discriminanten.}
\end{center}

Es sei
\weq{1}{f(\Theta) = \Theta^n + a_1\Theta^{n-1} + \cdots + a_n = 0}
die irreducible Gleichung $n^{\text{ten}}$ Grades mit rationalen Co\"efficienten $a_1, a_2, \ldots, a_n$, die uns einen algebraischen K\"orper $\Omega = R(\Theta)$ definirt. Der K\"orper $\Omega$ ist dann der Inbegriff aller Zahlen der Form
\weq{2}{\omega = h_1 + h_2\Theta + h_3\Theta^2 + \cdots + h_n\Theta^{n-1},}
worin $h_1, h_2, \ldots, h_n$ \gesperrt{rationale Zahlen} sind. Setzen wir aber f\"ur $h_1, h_2, \ldots, h_n$ rationale Functionale, so erhalten wir aus (2) alle \gesperrt{Functionale} des K\"orpers $\Omega$.

Betrachten wir ein beliebiges System von $n$ Zahlen
\weq{3}{\omega_r = h_{1,r} + h_{2,r}\Theta + \ldots + h_{n,r}\Theta^{n-1}, \quad r = 1,2,\ldots,n}
unter der Voraussetzung, dass die Determinante
\weq{4}{H = \Sigma \pm h_{1,1}h_{2,2}\ldots h_{n,n}}
von Null verschieden ist, so kann, wegen der Irreducibilit\"at von $f$, eine Gleichung der Form
\weq{5}{k_1\omega_1 + k_2\omega_2 + \cdots + k_n\omega_n = 0,}
in der $k_1, k_2, \ldots, k_n$ rationale Zahlen sind, nur dann bestehen, wenn diese Co\"efficienten alle Null sind, und dies gilt auch dann noch, wenn in der Gleichung (5) f\"ur die Co\"efficienten \gesperrt{rationale Functionale} zugelassen werden.
